%=========================================================================
\chapter{Code Generation}
\label{chap:codegen}
\renewcommand{\sectiontitle}{}
%=========================================================================

While all the previous chapters of this guide described how to construct, use and manipulate objects of the various classes forming the \gls{WIR} \gls{IR}, another central aspect of an \gls{IR} is to generate actual code from it for further processing by, e.g., assemblers or linkers.
\cref{sec:streamio} thus puts a focus on the generic mechanisms of \gls{WIR} to emit its code in a processor-independent and easily comprehensible format.
Manipulators that allow to increase or decrease the level of detail of such code output are presented in \cref{sec:streammanip}.
\cref{sec:asmcode} discusses how to adapt this code output such that syntactically correct assembly output for the processor architectures presented in \crefrange{ssec:tcmodel}{ssec:armmodel} can be obtained.
The \gls{WIR} system model features aspects of a global memory layout and allocation (cf. \cref{sec:memlayout}) which are taken up in \cref{sec:ldscripts} that describes the generation of tailored linker scripts for a given memory allocation.


%-------------------------------------------------------------------------
\section{Generic Stream Output}
\label{sec:streamio}
%-------------------------------------------------------------------------

Generating code from a given \gls{WIR} structure is as easy as a \texttt{"Hello world"} exercise in C++.
This is because all classes forming the \gls{WIR} software model (cf. \cref{chap:swmodel}) are equipped with \texttt{<{}<} stream operators that allow to dump the corresponding objects to an output stream in a highly flexible and configurable way.
More precisely, these operators exist for the following classes: \texttt{WIR\_System}, \texttt{WIR\_CompilationUnit}, \texttt{WIR\_Function}, \texttt{WIR\_BasicBlock}, \texttt{WIR\_Instruction}, \texttt{WIR\_Operation}, \texttt{WIR\_Parameter} and all derived parameter classes, and \texttt{WIR\_Data}.
They are all declared in analogy to \cref{lst:bbio} showing, e.g., the \gls{API} of class \texttt{WIR\_BasicBlock}.

\begin{lstlisting}[style=code,float=h,caption={Generic Stream Output of Class \texttt{WIR\_BasicBlock}},label=lst:bbio]
class WIR_BasicBlock
{
  public:

    [...]

    // Dump a basic block to an output stream.
    friend ästd::ostreamä æ&æ operator æ<<æ @(@ ästd::ostreamä æ&æos, const WIR_BasicBlock æ&æb @)@;Ü\label{lst:bbio1}Ü
};
\end{lstlisting}

An example program showing the use of such stream output and of various other concepts introduced in \cref{chap:swmodel,chap:hwmodel} is shown in \cref{lst:iotest}.
This example makes use of the RV32IMC processor model from \cref{ssec:rv32model} (cf. \cref{lst:iotest2}).
As already introduced in \cref{chap:build}, the \gls{WIR} library is initialized first in \cref{lst:iotest4}.
Hereafter, an RV32IMC-based system is created (\cref{lst:iotest5}) and its (one and only) processor core is obtained in \cref{lst:iotest6}.
\crefrange{lst:iotest7}{lst:iotest12} generate two compilation units containing one function each.
In the first function \texttt{foo} of compilation unit \texttt{file1.s}, three RISC-V virtual registers are created (\crefrange{lst:iotest13}{lst:iotest15}) where two of them are precolored with the processor's physical registers \texttt{x0} and \texttt{x1}, resp. (\cref{lst:iotest16,lst:iotest17}).
Next, one basic block containing one instruction is added to this function \texttt{foo} (\cref{lst:iotest18}), and this instruction contains two operations (\cref{lst:iotest19,lst:iotest22}).

\begin{lstlisting}[style=code,float=!ht,caption={Simple WIR Program Showing Stream Output},label=lst:iotest]
// Include WIR headers.
^#include^ §<wir/wir.h>§Ü\label{lst:iotest1}Ü
^#include^ §<arch/riscv/rv32imc.h>§Ü\label{lst:iotest2}Ü

using namespace WIR;Ü\label{lst:iotest3}Ü

@int@ main@( void )@
{
  // Initialize WIR.
  WIR_Init@()@;Ü\label{lst:iotest4}Ü

  // Create RISC-V WIR system and get the system's RV32IMC processor core.
  WIR_System sys@(@ "rudolv.sys", WIR_TaskManager() @)@;Ü\label{lst:iotest5}Ü
  auto æ&æp = sys.getComponentsæ<æRV32IMCæ>æ@()@.begin@()@æ->æget@()@;Ü\label{lst:iotest6}Ü

  // Create two compilation units with one function each.
  auto æ&æc1 = sys.pushBackCompilationUnit@(@ WIR_Function@(@ "foo" @) )@;Ü\label{lst:iotest7}Ü
  c1.setName@(@ "file1.s" @)@;Ü\label{lst:iotest8}Ü
  auto æ&æc2 = sys.pushBackCompilationUnit@(@ WIR_Function@(@ "bar" @) )@;Ü\label{lst:iotest9}Ü
  c2.setName@(@ "file2.s" @)@;Ü\label{lst:iotest10}Ü

  auto æ&æf1 = c1.getFunctions@()@.front@()@.get@()@;Ü\label{lst:iotest11}Ü
  auto æ&æf2 = c2.getFunctions@()@.front@()@.get@()@;Ü\label{lst:iotest12}Ü

  // Create various WIR objects in first compilation unit.
  auto æ&ær1 = f1.pushBackVirtualRegister@(@ WIR_VirtualRegister@(@ RV32I::RegisterType::reg @) )@;Ü\label{lst:iotest13}Ü
  auto æ&ær2 = f1.pushBackVirtualRegister@(@ WIR_VirtualRegister@(@ RV32I::RegisterType::reg @) )@;Ü\label{lst:iotest14}Ü
  auto æ&ær3 = f1.pushBackVirtualRegister@(@ WIR_VirtualRegister@(@ RV32I::RegisterType::reg @) )@;Ü\label{lst:iotest15}Ü
  f1.insertPrecolor@(@ r2, p.x0@() )@;Ü\label{lst:iotest16}Ü
  f1.insertPrecolor@(@ r3, p.x1@() )@;Ü\label{lst:iotest17}Ü

  f1.pushBackBasicBlock@(@ æ{æ WIR_Instruction@(@Ü\label{lst:iotest18}Ü
    WIR_Operation@(@ RV32I::OpCode::ADDI, RV32I::OperationFormat::RRC12_1,Ü\label{lst:iotest19}Ü
      WIR_RegisterParameter@(@ r1, WIR_Usage::def @)@, WIR_RegisterParameter@(@ r2, WIR_Usage::use @)@,Ü\label{lst:iotest20}Ü
      RV_Const12_Signed@(@ §42§ @) )@,Ü\label{lst:iotest21}Ü
    WIR_Operation@(@ RV32I::OpCode::JAL, RV32I::OperationFormat::RL_1,Ü\label{lst:iotest22}Ü
      WIR_RegisterParameter@(@ r3, WIR_Usage::def @)@, WIR_LabelParameter@(@ f2 @) ) )@ æ}æ @)@;Ü\label{lst:iotest23}Ü

  // Create various WIR objects in second compilation unit.
  auto æ&ær4 = f2.pushBackVirtualRegister@(@ WIR_VirtualRegister@(@ RV32I::RegisterType::reg @) )@;Ü\label{lst:iotest24}Ü
  auto æ&æd = c2.pushBackData@(@ æ{æ WIR_Data@(@ "dataObject" @)@ æ}æ @)@;Ü\label{lst:iotest25}Ü
  d.setSize@(@ §4§ @)@;Ü\label{lst:iotest26}Ü
  d.pushBackInitData@(@ WIR_DataInit@(@ WIR_DataInitType::iword, "4711" @) )@;Ü\label{lst:iotest27}Ü

  f2.pushBackBasicBlock@(@Ü\label{lst:iotest28}Ü
    æ{æ WIR_Instruction@(@ æ{æ RV32I::OpCode::LUI, RV32I::OperationFormat::RL_2,Ü\label{lst:iotest29}Ü
        WIR_RegisterParameter@(@ r4, WIR_Usage::def @)@, WIR_LabelParameter@(@ d @)@ æ}æ @)@,Ü\label{lst:iotest30}Ü
      WIR_Instruction@(@ æ{æ RV32I::OpCode::LW, RV32I::OperationFormat::RLR_1,Ü\label{lst:iotest31}Ü
         WIR_RegisterParameter@(@ p.x10@()@, WIR_Usage::def @)@, WIR_LabelParameter@(@ d @)@,Ü\label{lst:iotest32}Ü
         WIR_RegisterParameter@(@ r4, WIR_Usage::use @)@ æ}æ @)@,Ü\label{lst:iotest33}Ü
      WIR_Instruction@(@ æ{æ RV32I::OpCode::JALR, RV32I::OperationFormat::RRC12_1,Ü\label{lst:iotest34}Ü
         WIR_RegisterParameter@(@ p.x0@()@, WIR_Usage::def @)@, WIR_RegisterParameter@(@ p.x1@()@, WIR_Usage::use @)@,Ü\label{lst:iotest35}Ü
         RV_Const12_Signed@(@ §0§ @)@, WIR_RegisterParameter@(@ p.x10@()@, WIR_Usage::use, true @)@ æ}æ @)@ æ}æ @)@;Ü\label{lst:iotest36}Ü

  // Generate code for the current WIR system.
  ästd::coutä æ<<æ sys æ<<æ ästd::endlä;Ü\label{lst:iotest37}Ü

  return@(@ §0§ @)@;
}
\end{lstlisting}
\clearpage

In a similar fashion, another virtual register is created within the second function \texttt{bar} of compilation unit \texttt{file2.s} (\cref{lst:iotest24}).
\crefrange{lst:iotest25}{lst:iotest27} show the construction of a data object inside \texttt{file2.s} that is initialized with \texttt{4711}.
Function \texttt{bar} itself obtains one basic block (\cref{lst:iotest28}) featuring three instructions where each instruction consists of exactly one operation each (\cref{lst:iotest29,lst:iotest31,lst:iotest34}).

Finally, the entire code structure accommodated by the \gls{WIR} system is simply printed out to \texttt{cout} in \cref{lst:iotest37}.

Compiling and executing this example program from \cref{lst:iotest} according to \cref{chap:build} yields the following \gls{WIR} output:

\begin{lstlisting}[style=commandline,float=h]
FILE file1.sÜ\label{lst:wiroutput1}Ü
{Ü\label{lst:wiroutput2}Ü

SECTION .textÜ\label{lst:wiroutput3}Ü
{Ü\label{lst:wiroutput4}Ü

FUNCTION fooÜ\label{lst:wiroutput5}Ü
{Ü\label{lst:wiroutput6}Ü
foo:Ü\label{lst:wiroutput7}Ü
        addi RRC12_1 x_5733, x_5735, 42;Ü\label{lst:wiroutput8}Ü
     Ü||Ü jal RL_1 x_5737, bar;Ü\label{lst:wiroutput9}Ü
};Ü\label{lst:wiroutput10}Ü

};Ü\label{lst:wiroutput11}Ü

};Ü\label{lst:wiroutput12}Ü


FILE file2.sÜ\label{lst:wiroutput13}Ü
{Ü\label{lst:wiroutput14}Ü

SECTION .dataÜ\label{lst:wiroutput15}Ü
{Ü\label{lst:wiroutput16}Ü

DATA dataObjectÜ\label{lst:wiroutput17}Ü
{Ü\label{lst:wiroutput18}Ü
  4711Ü\label{lst:wiroutput19}Ü
};Ü\label{lst:wiroutput20}Ü

};Ü\label{lst:wiroutput21}Ü

SECTION .textÜ\label{lst:wiroutput22}Ü
{Ü\label{lst:wiroutput23}Ü

FUNCTION barÜ\label{lst:wiroutput24}Ü
{Ü\label{lst:wiroutput25}Ü
bar:Ü\label{lst:wiroutput26}Ü
        lui RL_2 x_5758, dataObject;Ü\label{lst:wiroutput27}Ü
        lw RLR_1 x10, dataObject, x_5758;Ü\label{lst:wiroutput28}Ü
        jalr RRC12_1 x0, x1, 0;Ü\label{lst:wiroutput29}Ü
};Ü\label{lst:wiroutput30}Ü

};Ü\label{lst:wiroutput31}Ü

};Ü\label{lst:wiroutput32}Ü
\end{lstlisting}

This output is entirely generic in that sense that it is unaware of any RISC-V specific code conventions.
As can be seen from, e.g., \cref{lst:wiroutput2,lst:wiroutput4,lst:wiroutput6} and \crefrange{lst:wiroutput10}{lst:wiroutput12}, such generic \gls{WIR} output consists of a dump of nested, hierarchical scopes that are all enclosed in curly brackets \texttt{\{\,...\,\}}.
Each such scope begins with a keyword specifying the scope's nature, e.g., \texttt{FILE}, \texttt{SECTION}, \texttt{FUNCTION} or \texttt{DATA}, followed by the (optional) name of the corresponding component.

The output of a \gls{WIR} function consists of the printout of all basic blocks of that function.
Each basic block is printed out beginning with its name followed by a colon (cf. \cref{lst:wiroutput7,lst:wiroutput26}).
Line by line, all instructions contained by a basic block are emitted next (e.g., \crefrange{lst:wiroutput27}{lst:wiroutput29}).
For each instruction, its operations are dumped again line by line, where the operation's output starts (indented by 8 whitespaces) with its mnemonic string followed by the operation format's name.
Finally, all explicit parameters of an operation are printed out as comma-separated list.
The output of an operation ends with a semicolon (\cref{lst:wiroutput28}).
For instructions containing several operations, the dump of all operations except the very first one always begins with a ``parallel'' operator \texttt{||} as can be seen in \cref{lst:wiroutput9}.
According to \cref{sec:instructions}, all operations of an instruction are supposed to be executed in parallel, which is indicated in the \gls{WIR} output this way.


%-------------------------------------------------------------------------
\section{Stream Manipulators}
\label{sec:streammanip}
%-------------------------------------------------------------------------

The generic \gls{WIR} output from the previous section reveals that certain code properties that are specified in the example program from \cref{lst:iotest} are not printed out by default:
\begin{itemize}
  \item Function \texttt{foo} contains virtual registers \texttt{x\_5733}, \texttt{x\_5735} and \texttt{x\_5737} (\crefrange{lst:iotest13}{lst:iotest15}), and function \texttt{bar} features virtual register \texttt{x\_5758} (\cref{lst:iotest24});
  \item Virtual registers \texttt{x\_5735} and \texttt{x\_5737} are precolored with certain physical registers (\cref{lst:iotest16,lst:iotest17});
  \item Various register parameters of the shown operations are defined (e.g., \cref{lst:iotest23}), while others are used (e.g., \cref{lst:iotest33});
  \item Implicit register parameters such as \texttt{p.x10} in \cref{lst:iotest36} are not printed at all in the default \gls{WIR} output.
\end{itemize}
This default behavior is motivated by the objective to keep the output as simple as possible and to primarily focus on the \gls{WIR} software model's code hierarchy.
However, the output of additional details can selectively be turned on or off by using the C++ stream manipulators declared in file \texttt{WIR/wir/wirio.h}.
A stream manipulator is a simple function that receives an output stream as argument, modifies it internally so that different output will be created, and returns a reference to the modified stream.
\cref{lst:iodefuse1} of \cref{lst:iodefuse} shows a sample declaration of the \texttt{defuse} stream manipulator.
This manipulator enables the output of additional def/use properties for register parameters.
In analogy to the standard C++ \gls{API}, each \gls{WIR} manipulator activating a certain output property comes with a manipulator to deactivate this behavior again.
The name of such manipulators always begins with ``\texttt{no}'' which can be seen from the example \texttt{nodefuse} from \cref{lst:iodefuse2}.

\begin{lstlisting}[style=code,float=h,caption={Declaration of \texttt{defuse} and \texttt{nodefuse} Stream Manipulators in file \texttt{wirio.h}},label=lst:iodefuse]
// Switch the output of def/use characteristics for register parameters on.
ästd::ostreamä æ&ædefuse@(@ ästd::ostreamä æ&æos @)@;Ü\label{lst:iodefuse1}Ü

// Switch the output of def/use characteristics for register parameters off.
ästd::ostreamä æ&ænodefuse@(@ ästd::ostreamä æ&æos @)@;Ü\label{lst:iodefuse2}Ü
\end{lstlisting}

In order to control the level of detail of generic \gls{WIR} output, the manipulators shown in \cref{tab:iomanip} are declared in file \texttt{wirio.h}.
They allow to address all the abovementioned issues.
In addition to that, the output of jump targets that are explicitly annotated for jump operations (cf. \cref{lst:opjmptgt}) can also be controlled.
Finally, the output of potentially existing comments and source file information can also be activated or deactivated.

When changing \cref{lst:iotest37} from \cref{lst:iotest} to

\begin{lstlisting}[style=code,float=h,firstnumber=55]
  // Generate code for the current WIR system.
  ästd::coutä æ<<æ defuse æ<<æ functionregisters æ<<æ implicitparams æ<<æ precolors æ<<æ sys æ<<æ ästd::endlä;
\end{lstlisting}

\begin{table}[t]
  \begin{center}
    \caption{Generic Stream Manipulators from file \texttt{wirio.h}}
    \label{tab:iomanip}
    \begin{tabular}{|l|p{9.2cm}|}
      \hline
      \textbf{Manipulators} & \textbf{Enabled/Disabled Output Property} \\
      \hhline{|=|=|}
      \texttt{comment}, \texttt{nocomment} & Comments attached to \gls{WIR} objects \\
      \hline
      \texttt{defuse}, \texttt{nodefuse} & Def/use characteristics of register parameters \\
      \hline
      \texttt{fileinfo}, \texttt{nofileinfo} & Source code file information attached to \gls{WIR} objects \\
      \hline
      \texttt{functionregisters}, \texttt{nofunctionregisters} & Virtual registers owned by functions \\
      \hline
      \texttt{implicitparams}, \texttt{noimplicitparams} & Implicit parameters of operations \\
      \hline
      \texttt{jumptargets}, \texttt{nojumptargets} & Explicitly specified targets of jump operations \\
      \hline
      \texttt{precolors}, \texttt{noprecolors} & Precoloring information of functions \\
      \hline
    \end{tabular}
  \end{center}
\end{table}

the produced \gls{WIR} output for the two functions \texttt{foo} and \texttt{bar} from the previous example program changes to:

\begin{lstlisting}[style=commandline,float=!h]
[...]

FUNCTION fooÜ\label{lst:wirmanip1}Ü
{Ü\label{lst:wirmanip2}Ü
  {Ü\label{lst:wirmanip3}Ü
    registers: x_5733 x_5735 x_5737;Ü\label{lst:wirmanip4}Ü
    precolors: x_5735=x0 x_5737=x1;Ü\label{lst:wirmanip5}Ü
  }Ü\label{lst:wirmanip6}Ü
foo:Ü\label{lst:wirmanip7}Ü
        addi RRC12_1 x_5733 (def, ID = 5747), x_5735 (use, ID = 5748), 42;Ü\label{lst:wirmanip8}Ü
     Ü||Ü jal RL_1 x_5737 (def, ID = 5741), bar;Ü\label{lst:wirmanip9}Ü
};Ü\label{lst:wirmanip10}Ü

[...]

FUNCTION barÜ\label{lst:wirmanip11}Ü
{Ü\label{lst:wirmanip12}Ü
  {Ü\label{lst:wirmanip13}Ü
    registers: x_5755;Ü\label{lst:wirmanip14}Ü
    precolors: ;Ü\label{lst:wirmanip15}Ü
  }Ü\label{lst:wirmanip16}Ü
bar:Ü\label{lst:wirmanip17}Ü
        lui RL_2 x_5755 (def, ID = 5762), dataObject;Ü\label{lst:wirmanip18}Ü
        lw RLR_1 x10 (def, ID = 5769), dataObject, x_5755 (use, ID = 5771);Ü\label{lst:wirmanip19}Ü
        jalr RRC12_1 x0 (def, ID = 5778), x1 (use, ID = 5779), 0, x10 (use, ID = 5781);Ü\label{lst:wirmanip20}Ü
};Ü\label{lst:wirmanip21}Ü

[...]
\end{lstlisting}

As can be seen from \crefrange{lst:wirmanip3}{lst:wirmanip6} and \crefrange{lst:wirmanip13}{lst:wirmanip16}, each function now begins with another scope containing information about the function's registers and its specified precolorings.
This information starts with the keywords \texttt{registers:} and \texttt{precolors:}, resp., followed by a whitespace-separated list of registers or precolorings.
These lines of output always end with a semicolon.
\cref{lst:wirmanip8,lst:wirmanip9,lst:wirmanip18,lst:wirmanip19,lst:wirmanip20} furthermore show that the output for each register parameter now is extended by parentheses \texttt{(\,...\,)} containing the information whether a parameter is defined, used, or both, plus the parameter's unique ID.
A closer look at \cref{lst:wirmanip20} also shows that the previously omitted implicit parameter is now printed, too.

\begin{lstlisting}[style=code,float=t,caption={Constructors and Methods for Comment Text Handling of Class \texttt{WIR\_Comment}},label=lst:commentconstruction]
// Default constructor creating a single-line comment from a copied string.
explicit WIR_Comment::WIR_Comment@(@ const ästd::stringä æ&æs @)@;Ü\label{lst:commentconstruction1}Ü

// Default constructor creating a single-line comment from a moved string.
explicit WIR_Comment::WIR_Comment@(@ ästd::stringä æ&&æs @)@;Ü\label{lst:commentconstruction2}Ü

// Copy constructor.
WIR_Comment::WIR_Comment@(@ const WIR_Comment æ&æc @)@;Ü\label{lst:commentconstruction3}Ü

// Move constructor.
WIR_Comment::WIR_Comment@(@ WIR_Comment æ&&æc @)@;Ü\label{lst:commentconstruction4}Ü

// Set comment text from a copied string.
@void@ WIR_Comment::setText@(@ const ästd::stringä æ&æs @)@;Ü\label{lst:commentconstruction5}Ü

// Set comment text from a moved string.
@void@ WIR_Comment::setText@(@ ästd::stringä æ&&æs @)@;Ü\label{lst:commentconstruction6}Ü

// Get a comment's text.
ästd::stringä WIR_Comment::getText@( void )@ const;Ü\label{lst:commentconstruction7}Ü
\end{lstlisting}

\begin{lstlisting}[style=code,float=!b,caption={Constructors and Methods for File Location Handling of Class \texttt{WIR\_FileInfo}},label=lst:fileinfoconstruction]
// Default constructor creating file information from a copied string.
WIR_FileInfo::WIR_FileInfo@(@ const ästd::stringä æ&æs, äunsigned long longä l @)@;Ü\label{lst:fileinfoconstruction1}Ü

// Default constructor creating file information from a moved string.
WIR_FileInfo::WIR_FileInfo@(@ ästd::stringä æ&&æs, äunsigned long longä l @)@;Ü\label{lst:fileinfoconstruction2}Ü

// Copy constructor.
WIR_FileInfo::WIR_FileInfo@(@ const WIR_FileInfo æ&æf @)@;Ü\label{lst:fileinfoconstruction3}Ü

// Move constructor.
WIR_FileInfo::WIR_FileInfo@(@ WIR_FileInfo æ&&æf @)@;Ü\label{lst:fileinfoconstruction4}Ü

// Set file name from a copied string.
@void@ WIR_FileInfo::setFileName@(@ const ästd::stringä æ&æs @)@;Ü\label{lst:fileinfoconstruction5}Ü

// Set file name from a moved string.
@void@ WIR_FileInfo::setFileName@(@ ästd::stringä æ&&æs @)@;Ü\label{lst:fileinfoconstruction6}Ü

// Get the file name.
ästd::stringä WIR_FileInfo::getFileName@( void )@ const;Ü\label{lst:fileinfoconstruction7}Ü

// Set line number.
@void@ WIR_FileInfo::setLineNumber@(@ äunsigned long longä l @)@;Ü\label{lst:fileinfoconstruction8}Ü

// Get line number.
äunsigned long longä WIR_FileInfo::getLineNumber@( void )@ const;Ü\label{lst:fileinfoconstruction9}Ü
\end{lstlisting}

\gls{WIR} uses so-called containers in order to attach arbitrary kinds of meta-data to all parts of its software model.
The discussion of this container \gls{API} is beyond the scope of this chapter, details will be provided later in \cref{sec:containers}.
However, the previously mentioned comments and source file information are very simple examples of such containers.

Class \texttt{WIR\_Comment} can be found in file \texttt{WIR/containers/wircomment.h}, and it represents single-line string comments.
The public interface of such comments (cf. \cref{lst:commentconstruction}) is extremely straightforward and simple so that a more detailed discussion is omitted here.
However, it is worthwhile to mention that comment containers represent single-line comments so that any newline character passed to class \texttt{WIR\_Comment} is removed.

File information objects intend to establish a link between a certain construct within the \gls{WIR} software model and a location in some source code from which the \gls{WIR} object has been generated.
Such code locations usually come in the form of pairs of file names and line numbers, and this is exactly what is represented by class \texttt{WIR\_FileInfo} from file \texttt{WIR/containers/wirfileinfo.h}, cf. \cref{lst:fileinfoconstruction}.
Both a file name and a line number can be specified during construction or can be set later via the \texttt{set} methods.
As above, file names are assumed to be single-line strings so that line breaks are removed.

Coming back to the example program from \cref{lst:iotest}, such comments and file information can be added and printed out, e.g., as follows:

\begin{lstlisting}[style=code,float=h,firstnumber=55]
  f1.insertContainer( WIR_FileInfo( "file1.c", 37 ) );Ü\label{lst:fileinfo1}Ü
  f2.begin()->get().insertContainer( WIR_Comment( "A commented WIR basic block." ) );Ü\label{lst:fileinfo2}Ü

  // Generate code for the current WIR system.
  ästd::coutä æ<<æ comment æ<<æ fileinfo æ<<æ sys æ<<æ ästd::endlä;Ü\label{lst:fileinfo3}Ü
\end{lstlisting}

\cref{lst:fileinfo1} attaches a file location to function \texttt{foo}, referring to some C file \texttt{file1.c} and line number 37.
The first basic block of function \texttt{bar} is extended by a comment in \cref{lst:fileinfo2}.
The output finally resulting from \cref{lst:fileinfo3} is:

\begin{lstlisting}[style=commandline,float=!h]
[...]

Ü\textit{\color{gray}; file "file1.c", line 37\label{lst:comment1}}Ü
FUNCTION foo
{
foo:
        addi RRC12_1 x_5733, x_5735, 42;
     Ü||Ü jal RL_1 x_5737, bar;
};

[...]

FUNCTION bar
{
Ü\textit{\color{gray}; A commented WIR basic block.\label{lst:comment2}}Ü
bar:
        lui RL_2 x_5755, dataObject;
        lw RLR_1 x10, dataObject, x_5755;
        jalr RRC12_1 x0, x1, 0;
};

[...]
\end{lstlisting}

As can be seen, both the file information (cf. \cref{lst:comment1}) and the comment (\cref{lst:comment2}) are printed before their respective \gls{WIR} objects (here: function \texttt{foo} and basic block \texttt{bar}) are printed.
The output of both containers starts with a semicolon followed by the container's data in the rest of the line.


%-------------------------------------------------------------------------
\section{Processor-Specific Assembly Code Generation}
\label{sec:asmcode}
%-------------------------------------------------------------------------

The discussion of code generation in the previous sections of this chapter solely focussed on generic and processor-independent output.
This section thus shows how to produce processor-specific output that obeys the assembly-level code conventions of the architectures currently supported by \gls{WIR}.

From a user's perspective, generating processor-specific assembly code is as simple as already introduced in \cref{sec:streammanip} by using the stream manipulators summarized in \cref{tab:procmanip}.
They are named according to their respective processor families shown in the second column and aim to provide GNU assembler-compliant output.
As shown in the right column, they are declared inside the architectures' various sub-directories.
Since the MIPS SPIM architecture just serves as a demonstrator within \gls{WIR} and is currently not used for real code generation, the \texttt{mips} manipulator only produces rudimentary valid assembly output.
For the other manipulators \texttt{arm}, \texttt{riscv} and \texttt{tricore}, the \gls{WIR} unit tests (cf. \cref{chap:build}) ensure the generation of valid assembly.
The \texttt{wir} stream manipulator switches a stream back to the generic output described in the previous sections.

When revisiting the previous example, generating RISC-V code can be achieved easily by changing \cref{lst:fileinfo3} to:

\begin{lstlisting}[style=code,float=h,firstnumber=58]
  // Generate RISC-V code for the current WIR system.
  ästd::coutä æ<<æ riscv æ<<æ comment æ<<æ fileinfo æ<<æ sys æ<<æ ästd::endlä;Ü\label{lst:rvoutput}Ü
\end{lstlisting}

As a result, \gls{WIR} now produces the following assembly output:

\begin{lstlisting}[style=commandline,float=!h]
        .file           "file1.s"
        # Generated using wir release 2.0 (2024-03-24)
        # for RISC-V ISA RV32IMC

        .section        .text
        .balign         8
        # file "file1.c", line 37
        .type           foo, Ü@Üfunction
foo:
        addi            x_5733, x_5735/x0, 42
        jal             x_5737/x1, bar
        .size           foo, .-foo

        .file           "file2.s"
        # Generated using wir release 2.0 (2024-03-24)
        # for RISC-V ISA RV32IMC

        .section        .data, "aw"
        .balign         8
        .type           dataObject, Ü@Üobject
        .size           dataObject, 4
dataObject:
        .word           4711

        .section        .text
        .balign         8
        .type           bar, Ü@Üfunction
        # A commented WIR basic block.
bar:
        lui             x_5755, %hi(dataObject)Ü\label{lst:rvasm1}Ü
        lw              x10, %lo(dataObject)(x_5755)Ü\label{lst:rvasm2}Ü
        jalr            x0, x1, 0
        .size           bar, .-bar
\end{lstlisting}

\begin{table}[t]
  \begin{center}
    \caption{Processor-Specific Stream Manipulators}
    \label{tab:procmanip}
    \begin{tabular}{|l|l|p{5.1cm}|}
      \hline
      \textbf{Manipulator} & \textbf{Processor Families} & \textbf{File of Declaration} \\
      \hhline{|=|=|=|}
      \texttt{arm} & All models from ARMv4 to ARMv6 from \cref{ssec:armmodel} & \texttt{WIR/arch/arm/armio.h} \\
      \hline
      \texttt{mips} & MIPS SPIM model from \cref{ssec:mipsmodel} & \texttt{WIR/arch/generic/mipsio.h} \\
      \hline
      \texttt{riscv} & All models from RV32I to RV32IMFDC from \cref{ssec:rv32model} & \texttt{WIR/arch/riscv/rvio.h} \\
      \hline
      \texttt{tricore} & TC13 and TC131 models from \cref{ssec:tcmodel} & \texttt{WIR/arch/tricore/tcio.h} \\
      \hline
      \texttt{wir} & \gls{WIR} generic output from \cref{sec:streamio} & \texttt{WIR/wir/wirio.h} \\
      \hline
    \end{tabular}
  \end{center}
\end{table}

As can be seen, the various scopes of the previous, generic \gls{WIR} output representing, e.g., compilation units, sections and data objects are now printed as valid assembly directives.
The code of the functions themselves is also formatted in a proper manner where, e.g., the output of operation formats is omitted, or the specifiers \texttt{\%hi(...)} and \texttt{\%lo(...)} are used in \cref{lst:rvasm1,lst:rvasm2}.
The only aspect in the above example output that is not fully RISC-V assembly compliant is the use of virtual registers and their output together with their precolorings.
This, however, is not an issue of \gls{WIR}'s code generation mechanisms.
Instead, this is solely due to the use of virtual instead of physical registers in the running example.

For all classes of the \gls{WIR} software model, their respective \texttt{<{}<} stream operators invoke so-called dumper functions, where such a dumper performs the actual task of printing out a \gls{WIR} object.
File \texttt{WIR/wir/wirio.h} declares a couple of data types that represent function pointers for all these available dumper functions.
As \cref{lst:bbio} showed the example of the \texttt{<{}<} operator for basic blocks, \cref{lst:iodumper} now shows the declaration of the corresponding basic block dumper type.
Functions of this type thus expect a reference to an output stream as well as a constant reference to a basic block to be printed out.

\begin{lstlisting}[style=code,float=h,caption={Declaration of Basic Block Dumper Type in file \texttt{wirio.h}},label=lst:iodumper]
// Type WIR_BasicBlockDumper represents pointers to functions for processor-specific output of basic blocks.
using WIR_BasicBlockDumper = @void (@æ*æ@)(@ ästd::ostreamä æ&æ, const WIR_BasicBlock æ&æ @)@;Ü\label{lst:iodumper1}Ü
\end{lstlisting}

For all processor families listed in \cref{tab:procmanip}, there now exist actual dumper functions that realize the processor-specific output of \gls{WIR} objects.
For example, the RISC-V-specific output of basic blocks is realized by function \texttt{dumpRVBasicBlock} declared in file \texttt{WIR/arch/riscv/rvio.h} (cf. \cref{lst:rvdumper}).
As can be seen, this actual function is precisely of type \texttt{WIR\_BasicBlockDumper} as declared in \cref{lst:iodumper}.

\begin{lstlisting}[style=code,float=h,caption={Declaration of RISC-V-Specific Basic Block Dumper},label=lst:rvdumper]
// Dump a basic block to an output stream in a RISC-V-specific fashion.
@void@ dumpRVBasicBlock@(@ ästd::ostreamä æ&æos, const WIR_BasicBlock æ&æb @)@;Ü\label{lst:rvdumper1}Ü
\end{lstlisting}

The \gls{WIR} library now finally uses an internal map that returns the appropriate function pointers to these processor-specific dumpers, depending on the stream manipulator from \cref{tab:procmanip} that has been applied to an output stream before.
\cref{sssec:opformats} has already introduced that certain aspects of a processor model need to be registered centrally within \gls{WIR} during the initialization of the entire library.
This previous discussion of processor-specific initialization revolved around the registration of operation formats and opcodes, cf. \cref{lst:tcopfreg} on \cpageref{lst:tcopfreg}.
The last aspect of \gls{WIR} initialization now is the proper registration of the various processor-specific dumpers in this internal map.
The corresponding excerpt of the initialization for the RISC-V processor family is shown in \cref{lst:rvdumpreg}.
Here, all the existing dumpers are registered, and \cref{lst:rvdumpreg1} showcases the concrete example of the RISC-V-specific basic block dumper registration.

All registrations of dumpers from \crefrange{lst:rvdumpreg1}{lst:rvdumpreg2} use method \texttt{getProcessorTypeID} (cf. \cref{lst:proctypes} on \cpageref{lst:proctypes}) in order to associate the given dumpers with the unique constant ID of the currently initialized processor family.
The stream manipulators from \cref{tab:procmanip} now finally associate exactly this unique processor ID with an output stream so that all the \texttt{<{}<} operators of the \gls{WIR} software model can obtain the proper processor-specific dumpers from the abovementioned internal map.

\begin{lstlisting}[style=code,float=t,caption={Dumper Function Registration (here for Class \texttt{RV32I})},label=lst:rvdumpreg]
@void@ RV32I::init@( void )@
{
  [...]

  // Registration of RV32I dumpers.
  registerBasicBlockDumper( getProcessorTypeID(), dumpRVBasicBlock );Ü\label{lst:rvdumpreg1}Ü
  registerCompilationUnitDumper( getProcessorTypeID(), dumpRVCompilationUnit );
  registerDataDumper( getProcessorTypeID(), dumpRVData );
  registerDataSectionDumper( getProcessorTypeID(), dumpRVDataSection );
  registerFunctionDumper( getProcessorTypeID(), dumpRVFunction );
  registerOperationDumper( getProcessorTypeID(), dumpRVOperation );
  registerRegisterParameterDumper( getProcessorTypeID(), dumpRVRegisterParameter );
  registerCommentDumper( getProcessorTypeID(), dumpRVComment );
  registerFileInfoDumper( getProcessorTypeID(), dumpRVFileInfo );Ü\label{lst:rvdumpreg2}Ü

  // RV32I operation formats.

  [...]
};
\end{lstlisting}


%-------------------------------------------------------------------------
\section{Linker Script Generation}
\label{sec:ldscripts}
%-------------------------------------------------------------------------

One particular stream manipulator has not yet been introduced in the previous sections, as its purpose is not directly related to the generation of actual code from the \gls{WIR} software model.
As discussed earlier in \cref{sec:memlayout}, \gls{WIR} systems provide mechanisms for the mapping of logical sections to physical memory regions, and for the allocation of symbols to sections.
In other words, the allocation of code or data to a system's different memories can be determined at the system level.

\gls{WIR} is designed to be an assembly-level \gls{IR}.
In contrast, such memory allocation decisions are regularly enforced at the level of binary executable code by applying some linker.
The code generated using the concepts described in the previous sections of this chapter does not contain any provisions to enforce some memory allocation specified by a \gls{WIR} system.
For this purpose, the manipulator \texttt{ldscript} (cf. \cref{lst:ioldscript1} of \cref{lst:ioldscript}) can be used.
It enables the generation of a GNU linker-compliant script that in turn forces the linker to respect the memory allocations defined by the \gls{WIR} system.
This manipulator has only effect on the \texttt{<{}<} output operator of class \texttt{WIR\_System}, and the manipulator \texttt{noldscript} (\cref{lst:ioldscript2}) switches a stream back to the generation of regular code.

\begin{lstlisting}[style=code,float=h,caption={Declaration of \texttt{ldscript} and \texttt{noldscript} Stream Manipulators in file \texttt{wirio.h}},label=lst:ioldscript]
// Switch the output of a linker script for WIR systems on.
ästd::ostreamä æ&ældscript@(@ ästd::ostreamä æ&æos @)@;Ü\label{lst:ioldscript1}Ü

// Switch the output of a linker script for WIR systems off.
ästd::ostreamä æ&ænoldscript@(@ ästd::ostreamä æ&æos @)@;Ü\label{lst:ioldscript2}Ü
\end{lstlisting}


\cleardoubleoddpage
